% ============================================================
%  RAPPORT TECHNIQUE N°2 - Deep Learning
%  Algorithme de Rétropropagation & Réseaux Dense/Convolutif
%  Auteurs : Souaibou DIENG & Ibrahima FALL
%  ESP DAKAR - Master 2 SRT - 2025/2026
% ============================================================

\documentclass[12pt, a4paper]{article}

% ---------- Encodage & Langue ----------
\usepackage[utf8]{inputenc}
\usepackage[T1]{fontenc}
\usepackage[french]{babel}

% ---------- Mise en page ----------
\usepackage[top=2.5cm, bottom=2.5cm, left=3cm, right=2.5cm]{geometry}
\usepackage{setspace}
\onehalfspacing

% ---------- Polices ----------
\usepackage{lmodern}
\usepackage{microtype}

% ---------- Graphiques & Couleurs ----------
\usepackage{graphicx}
\usepackage{xcolor}
\usepackage{float}
\usepackage{tikz}
\usetikzlibrary{positioning, arrows.meta, shapes.geometric}

% Palette de couleurs
\definecolor{espblue}{RGB}{0, 70, 127}
\definecolor{espgold}{RGB}{204, 153, 0}
\definecolor{codegray}{RGB}{40, 44, 52}
\definecolor{codegreen}{RGB}{80, 200, 120}
\definecolor{codeorange}{RGB}{230, 126, 34}
\definecolor{codecomment}{RGB}{127, 140, 141}
\definecolor{codered}{RGB}{231, 76, 60}

% ---------- Tableaux ----------
\usepackage{booktabs}
\usepackage{tabularx}
\usepackage{array}
\usepackage{colortbl}

% ---------- Maths ----------
\usepackage{amsmath}
\usepackage{amssymb}
\usepackage{amsfonts}

% ---------- Code source ----------
\usepackage{listings}
\lstdefinestyle{pythonstyle}{
    backgroundcolor=\color{codegray},
    basicstyle=\ttfamily\footnotesize\color{white},
    keywordstyle=\color{codeorange}\bfseries,
    stringstyle=\color{codegreen},
    commentstyle=\color{codecomment}\itshape,
    numberstyle=\tiny\color{gray},
    numbers=left,
    numbersep=8pt,
    breaklines=true,
    breakatwhitespace=true,
    frame=single,
    framesep=5pt,
    rulecolor=\color{espblue},
    showstringspaces=false,
    tabsize=4,
    captionpos=b,
    literate=
        {é}{{\'{e}}}1 {è}{{\`{e}}}1 {ê}{{\^{e}}}1 {ë}{{\"e}}1
        {à}{{\`{a}}}1 {â}{{\^{a}}}1 {î}{{\^i}}1 {ô}{{\^o}}1
        {ù}{{\`{u}}}1 {û}{{\^{u}}}1 {ç}{{\c{c}}}1
}
\lstset{style=pythonstyle}

% ---------- Hyperliens ----------
\usepackage[hidelinks]{hyperref}
\hypersetup{
    pdftitle={Rapport TP2 - Rétropropagation \& Réseaux Dense/Convolutif},
    pdfauthor={Souaibou DIENG \& Ibrahima FALL},
}

% ---------- En-têtes et pieds de page ----------
\usepackage{fancyhdr}
\pagestyle{fancy}
\fancyhf{}
\fancyhead[L]{\small\textcolor{espblue}{\textit{Rétropropagation \& Réseaux Dense/Convolutif}}}
\fancyhead[R]{\small\textcolor{espblue}{\textit{Master 2 SRT -- ESP Dakar}}}
\fancyfoot[C]{\textcolor{espblue}{\thepage}}
\renewcommand{\headrulewidth}{0.4pt}
\renewcommand{\headrule}{\hbox to\headwidth{%
    \color{espblue}\leaders\hrule height \headrulewidth\hfill}}

% ---------- Boîtes ----------
\usepackage{tcolorbox}
\tcbuselibrary{skins, breakable}

\newtcolorbox{infobox}[1][]{
    enhanced, breakable,
    colback=espblue!5, colframe=espblue,
    boxrule=0.8pt, arc=4pt,
    fonttitle=\bfseries\color{white},
    colbacktitle=espblue, #1
}

\newtcolorbox{mathbox}{
    enhanced,
    colback=espgold!8, colframe=espgold,
    boxrule=0.8pt, arc=4pt,
}

\newtcolorbox{resultbox}{
    enhanced,
    colback=codegreen!8, colframe=codegreen!60!black,
    boxrule=0.8pt, arc=4pt,
}

\newtcolorbox{stepbox}[1]{
    enhanced, breakable,
    colback=espblue!4, colframe=espblue!35,
    boxrule=0.6pt, arc=3pt,
    title={\textbf{#1}},
    fonttitle=\small\bfseries\color{espblue},
    colbacktitle=espblue!12,
    coltitle=espblue,
}

\newtcolorbox{gradbox}{
    enhanced,
    colback=codered!5, colframe=codered!60,
    boxrule=0.6pt, arc=3pt,
}

% ---------- Sections personnalisées ----------
\usepackage{titlesec}
\titleformat{\section}
    {\normalfont\Large\bfseries\color{espblue}}
    {\thesection.}{10pt}{}
    [\vspace{2pt}{\color{espblue}\titlerule[0.6pt]}]
\titlespacing{\section}{0pt}{18pt}{10pt}

\titleformat{\subsection}
    {\normalfont\large\bfseries\color{espblue!80!black}}
    {\thesubsection}{8pt}{}
\titlespacing{\subsection}{0pt}{12pt}{6pt}

\titleformat{\paragraph}[runin]
    {\normalfont\normalsize\bfseries\color{espblue!70!black}}
    {}{0pt}{}[\ ---]

% ---------- Divers ----------
\usepackage{enumitem}
\usepackage{caption}
\captionsetup{font=small, labelfont=bf, labelsep=period}

% ============================================================
%  DÉBUT DU DOCUMENT
% ============================================================
\begin{document}

% ============================================================
%  PAGE DE GARDE
% ============================================================
\begin{titlepage}
    \centering
    \vspace*{0.3cm}

    {\color{espblue}\rule{\linewidth}{3pt}}\\[0.5cm]

    \begin{minipage}[c]{0.30\textwidth}
        \centering
        {\small\bfseries RÉPUBLIQUE DU SÉNÉGAL}\\[3pt]
        {\footnotesize Un Peuple – Un But – Une Foi}
    \end{minipage}
    \hfill
    \begin{minipage}[c]{0.30\textwidth}
        \centering
        \includegraphics[width=2.8cm]{logo_esp.png}
    \end{minipage}
    \hfill
    \begin{minipage}[c]{0.30\textwidth}
        \centering
        {\small\bfseries UNIVERSITÉ CHEIKH ANTA DIOP}\\[3pt]
        {\footnotesize DE DAKAR (UCAD)}
    \end{minipage}

    \vspace{0.5cm}
    {\color{espblue}\rule{\linewidth}{1pt}}\\[0.3cm]

    {\large\bfseries\color{espblue} ÉCOLE SUPÉRIEURE POLYTECHNIQUE DE DAKAR}\\[4pt]
    {\normalsize DÉPARTEMENT DE GÉNIE INFORMATIQUE}\\[2pt]
    {\normalsize\bfseries MASTER 2 -- SYSTÈMES, RÉSEAUX ET TÉLÉCOMMUNICATIONS}

    \vspace{0.4cm}
    {\color{espgold}\rule{0.5\linewidth}{2pt}}
    \vspace{0.3cm}

    {\large\color{espblue!70!black} Module :}\\[3pt]
    {\Large\bfseries Deep Learning / Traitement d'Images}

    \vspace{0.3cm}
    {\color{espgold}\rule{0.8\linewidth}{1.5pt}}
    \vspace{0.3cm}

    % Badge numéro de TP
    \begin{tcolorbox}[enhanced, width=4cm, colback=espblue, colframe=espblue,
        arc=20pt, boxrule=0pt, halign=center]
        {\large\bfseries\color{white} Compte Rendu TD}\\[2pt]
        {\LARGE\bfseries\color{espgold} N° 2}
    \end{tcolorbox}

    \vspace{0.3cm}

    {\LARGE\bfseries\color{espblue}
        Algorithme de Rétropropagation\\[5pt]
        et Réseaux Dense / Convolutif
    }

    \vspace{0.3cm}
    {\color{espgold}\rule{0.8\linewidth}{1.5pt}}
    \vspace{0.5cm}

    \Large\bfseries Réalisé par :\\[0.4cm]
    \begin{tabular}{cc}
        \Large\color{espblue} Souaibou \textsc{Dieng} &
        \Large\color{espblue} Ibrahima \textsc{Fall} \\
    \end{tabular}

    \vfill

    \begin{tabular}{rl}
        \textbf{Classe :}              & Master 2 SRT \\[3pt]
        \textbf{Encadrant :}           & Dr. Mouhamed Kaba \textsc{Keïta} \\[3pt]
        \textbf{Année universitaire :} & 2025 -- 2026 \\
    \end{tabular}

    \vspace{0.4cm}
    {\color{espblue}\rule{\linewidth}{3pt}}
\end{titlepage}

% ============================================================
%  TABLE DES MATIÈRES
% ============================================================
\tableofcontents
\thispagestyle{fancy}
\newpage

% ============================================================
%  INTRODUCTION
% ============================================================
\section*{Introduction}
\addcontentsline{toc}{section}{Introduction}

L'algorithme de \textbf{rétropropagation du gradient}
(\textit{Backpropagation}) constitue le moteur de l'apprentissage profond.
Il permet de mettre à jour les poids d'un réseau de neurones en calculant
l'influence de chaque paramètre sur l'erreur finale, grâce à la
\textbf{règle de dérivation en chaîne} (\textit{Chain Rule}).

Ce rapport détaille ce processus à travers deux exercices :
\begin{enumerate}[leftmargin=2cm, label=\textbf{\arabic{*}.}]
    \item Un \textbf{réseau dense (MLP 2-2-2)} avec des calculs numériques complets,
          passe avant et passe arrière.
    \item Une \textbf{étude théorique} sur la propagation des gradients dans les
          couches de convolution (CNN).
\end{enumerate}

% ============================================================
%  EXERCICE 1 : MLP 2-2-2
% ============================================================
\section{Étude d'un Réseau de Neurones Dense (MLP 2-2-2)}

On considère un réseau à deux entrées ($i_1, i_2$), une couche cachée
($h_1, h_2$) et deux sorties ($o_1, o_2$).

\subsection{Architecture et Paramètres Initiaux}

\begin{figure}[H]
    \centering
    \includegraphics[width=0.72\textwidth]{mlp222.png}
    \caption{Architecture du MLP 2-2-2 avec poids initiaux et biais.}
    \label{fig:mlp}
\end{figure}

\begin{table}[H]
    \centering
    \caption{Paramètres initiaux du réseau MLP 2-2-2}
    \label{tab:params}
    \renewcommand{\arraystretch}{1.4}
    \begin{tabularx}{0.9\textwidth}{lX}
        \toprule
        \rowcolor{espblue!15}
        \textbf{Paramètre} & \textbf{Valeur} \\
        \midrule
        Entrées           & $i_1 = 0.05$,\quad $i_2 = 0.10$ \\
        Poids couche 1    & $w_1=0.15,\ w_2=0.20,\ w_3=0.25,\ w_4=0.30$ \\
        Poids couche 2    & $w_5=0.40,\ w_6=0.45,\ w_7=0.50,\ w_8=0.55$ \\
        Biais             & $b_1 = 0.35$,\quad $b_2 = 0.60$ \\
        Cibles            & $t_{o1} = 0.01$,\quad $t_{o2} = 0.99$ \\
        Taux d'apprentissage & $\eta = 0.5$ \\
        \bottomrule
    \end{tabularx}
\end{table}

\subsection{Passe Avant (\textit{Forward Pass})}

Nous utilisons la \textbf{fonction d'activation sigmoïde} :

\begin{mathbox}
\begin{equation}
    \sigma(z) = \frac{1}{1 + e^{-z}}, \qquad
    \sigma'(z) = \sigma(z)\,\bigl(1 - \sigma(z)\bigr)
    \label{eq:sigmoid}
\end{equation}
\end{mathbox}

\begin{stepbox}{Étape 1 — Calcul de la couche cachée}
\begin{align*}
    net_{h1} &= w_1\,i_1 + w_3\,i_2 + b_1
               = 0.15(0.05) + 0.25(0.10) + 0.35 = 0.3775 \\
    out_{h1} &= \sigma(0.3775) \approx \mathbf{0.5933} \\[6pt]
    net_{h2} &= w_2\,i_1 + w_4\,i_2 + b_1
               = 0.20(0.05) + 0.30(0.10) + 0.35 = 0.3925 \\
    out_{h2} &= \sigma(0.3925) \approx \mathbf{0.5969}
\end{align*}
\end{stepbox}

\vspace{0.4cm}

\begin{stepbox}{Étape 2 — Calcul de la couche de sortie}
\begin{align*}
    net_{o1} &= w_5\,out_{h1} + w_7\,out_{h2} + b_2
               = 0.40(0.5933) + 0.50(0.5969) + 0.60 = 1.1059 \\
    out_{o1} &= \sigma(1.1059) \approx \mathbf{0.7514} \\[6pt]
    net_{o2} &= w_6\,out_{h1} + w_8\,out_{h2} + b_2
               = 0.45(0.5933) + 0.55(0.5969) + 0.60 = 1.2249 \\
    out_{o2} &= \sigma(1.2249) \approx \mathbf{0.7729}
\end{align*}
\end{stepbox}

\subsection{Calcul de l'Erreur Totale}

L'erreur est mesurée via l'\textbf{Erreur Quadratique Moyenne (MSE)} :

\begin{mathbox}
\begin{equation}
    E_{\text{total}} = \sum \frac{1}{2}(t - out)^2
    \label{eq:mse}
\end{equation}
\end{mathbox}

\begin{table}[H]
    \centering
    \caption{Décomposition de l'erreur totale}
    \label{tab:erreur}
    \renewcommand{\arraystretch}{1.4}
    \begin{tabular}{lcc}
        \toprule
        \rowcolor{espblue!15}
        \textbf{Sortie} & \textbf{Formule} & \textbf{Valeur} \\
        \midrule
        $E_{o1}$ & $\frac{1}{2}(0.01 - 0.7514)^2$ & $0.2748$ \\
        $E_{o2}$ & $\frac{1}{2}(0.99 - 0.7729)^2$ & $0.0236$ \\
        \midrule
        \rowcolor{codegreen!15}
        $E_{\text{total}}$ & $E_{o1} + E_{o2}$ & $\mathbf{0.2984}$ \\
        \bottomrule
    \end{tabular}
\end{table}

\subsection{Passe Arrière (\textit{Backward Pass}) et Mise à Jour}

La règle générale de mise à jour des poids est :

\begin{mathbox}
\begin{equation}
    w_{\text{new}} = w_{\text{old}} - \eta\,\frac{\partial E_{\text{total}}}{\partial w}
    \label{eq:update}
\end{equation}
\end{mathbox}

\begin{stepbox}{Mise à jour de $w_8$ — Couche de sortie}
Par la règle de la chaîne :
\[
    \frac{\partial E_{\text{total}}}{\partial w_8}
    = \underbrace{\frac{\partial E_{\text{total}}}{\partial out_{o2}}}_{-(t_{o2}-out_{o2})}
    \cdot \underbrace{\frac{\partial out_{o2}}{\partial net_{o2}}}_{out_{o2}(1-out_{o2})}
    \cdot \underbrace{\frac{\partial net_{o2}}{\partial w_8}}_{out_{h2}}
\]
\begin{align*}
    \frac{\partial E_{\text{total}}}{\partial out_{o2}}
        &= -(0.99 - 0.7729) = -0.2171 \\
    \frac{\partial out_{o2}}{\partial net_{o2}}
        &= 0.7729\,(1 - 0.7729) = 0.1755 \\
    \frac{\partial net_{o2}}{\partial w_8}
        &= out_{h2} = 0.5969 \\[4pt]
    \frac{\partial E_{\text{total}}}{\partial w_8}
        &\approx -0.0227 \\[4pt]
    w_8^{+} &= 0.55 - 0.5\,(-0.0227) = \mathbf{0.5613}
\end{align*}
\end{stepbox}

\vspace{0.4cm}

\begin{stepbox}{Mise à jour de $w_4$ — Couche cachée}
Le neurone $h_2$ contribue simultanément à $o_1$ et $o_2$, ce qui complexifie le calcul :
\[
    \frac{\partial E_{\text{total}}}{\partial out_{h2}}
    = \frac{\partial E_{o1}}{\partial out_{h2}}
    + \frac{\partial E_{o2}}{\partial out_{h2}}
\]
\begin{align*}
    \frac{\partial E_{\text{total}}}{\partial w_4}
        &\approx 0.0008 \\[4pt]
    w_4^{+} &= 0.30 - 0.5\,(0.0008) \approx \mathbf{0.2996}
\end{align*}
\end{stepbox}

\begin{resultbox}
\textbf{Synthèse :} Les nouveaux poids $w_8^+ = 0.5613$ et $w_4^+ = 0.2996$
illustrent que la rétropropagation ajuste chaque poids \textbf{proportionnellement
à sa contribution à l'erreur}, avec un pas contrôlé par $\eta$.
\end{resultbox}

% ============================================================
%  SECTION 3 : CODE PYTHON
% ============================================================
\newpage
\section{Validation par l'Implémentation Python}

Le script suivant (NumPy) vérifie les calculs manuels de la passe avant
et de la mise à jour de $w_8$ pour la première itération.

\begin{lstlisting}[language=Python, caption={Simulation du Backpropagation -- MLP 2-2-2}, label=lst:backprop]
import numpy as np

def sigmoid(x):
    return 1 / (1 + np.exp(-x))

def sigmoid_derivative(x):
    return x * (1 - x)

# --- Donnees initiales ---
i1, i2 = 0.05, 0.10
t1, t2 = 0.01, 0.99
eta    = 0.5

# --- Poids et biais ---
w1, w2, w3, w4 = 0.15, 0.20, 0.25, 0.30
w5, w6, w7, w8 = 0.40, 0.45, 0.50, 0.55
b1, b2         = 0.35, 0.60

# === PASSE AVANT ===
h1 = sigmoid(w1*i1 + w3*i2 + b1)   # out_h1
h2 = sigmoid(w2*i1 + w4*i2 + b1)   # out_h2
o1 = sigmoid(w5*h1 + w7*h2 + b2)   # out_o1
o2 = sigmoid(w6*h1 + w8*h2 + b2)   # out_o2

print("=== Passe Avant ===")
print(f"out_h1 = {h1:.4f}  |  out_h2 = {h2:.4f}")
print(f"out_o1 = {o1:.4f}  |  out_o2 = {o2:.4f}")
E_total = 0.5*(t1-o1)**2 + 0.5*(t2-o2)**2
print(f"E_total = {E_total:.4f}")

# === PASSE ARRIERE (w8) ===
dE_dout_o2  = -(t2 - o2)
dout_dnet   = sigmoid_derivative(o2)
dnet_dw8    = h2

grad_w8 = dE_dout_o2 * dout_dnet * dnet_dw8
w8_new  = w8 - eta * grad_w8

print("\n=== Passe Arriere ===")
print(f"Gradient w8 = {grad_w8:.4f}")
print(f"w8 : {w8:.4f}  ->  {w8_new:.4f}")
\end{lstlisting}

% ============================================================
%  EXERCICE 2 : CONVOLUTION
% ============================================================
\section{Convolution et Rétropropagation}

Cette partie analyse comment les gradients circulent au sein d'une
\textbf{couche de convolution} (CNN).

\subsection{Définition de la Sortie d'une Couche de Convolution}

Un élément de la carte de sortie $O$ est calculé par :

\begin{mathbox}
\begin{equation}
    O_{ij} = \sum_{p,\,q} X_{i+p,\,j+q} \cdot F_{pq}
    \label{eq:conv}
\end{equation}
où $X$ est l'image d'entrée et $F$ le filtre (noyau de convolution).
\end{mathbox}

\subsection{Gradients de la Perte $L$}

\begin{stepbox}{Gradient par rapport au filtre $F$}
\begin{equation}
    \frac{\partial L}{\partial F_{pq}}
    = \sum_{i,j} \frac{\partial L}{\partial O_{ij}} \cdot X_{i+p,\,j+q}
    \label{eq:grad_F}
\end{equation}
Le gradient par rapport au filtre est une \textbf{convolution} entre l'entrée $X$
et le gradient de la sortie $\delta O$.
\end{stepbox}

\vspace{0.4cm}

\begin{stepbox}{Gradient par rapport à l'entrée $X$}
La contribution de chaque élément $X_{mn}$ à la sortie dépend de sa position
sous le filtre. Le gradient s'écrit :
\begin{equation}
    \frac{\partial L}{\partial X_{mn}}
    = \sum_{i,j} \frac{\partial L}{\partial O_{ij}} \cdot F_{m-i,\,n-j}
    \label{eq:grad_X}
\end{equation}
Ce calcul correspond à une \textbf{corrélation croisée} entre le gradient de la
sortie et le filtre \textbf{inversé} (rotation de 180°).
\end{stepbox}

\begin{infobox}[title={Propriété fondamentale de la rétropropagation en CNN}]
Les deux opérations clés du backward pass en convolution sont :
\begin{itemize}[noitemsep]
    \item \textbf{Mise à jour du filtre} : convolution entre $X$ et $\delta O$
    \item \textbf{Propagation du gradient} : corrélation entre $\delta O$ et $F^{\text{rot180}}$
\end{itemize}
Ces propriétés permettent de conserver la \textbf{structure spatiale} des images
lors de l'apprentissage, ce qui est la force fondamentale des CNN.
\end{infobox}

% ============================================================
%  CONCLUSION
% ============================================================
\section*{Conclusion}
\addcontentsline{toc}{section}{Conclusion}

Ce rapport a démontré la puissance et l'élégance de la \textbf{règle de
dérivation en chaîne} (\textit{Chain Rule}) appliquée aux réseaux de neurones.

\begin{resultbox}
\textbf{Deux leçons essentielles :}
\begin{enumerate}[noitemsep, leftmargin=1.5cm]
    \item Dans un \textbf{réseau dense}, chaque poids est ajusté proportionnellement
          à sa contribution à l'erreur, remontée couche par couche.
    \item Dans une \textbf{couche convolutive}, l'apprentissage repose sur des
          opérations de corrélation qui préservent les propriétés spatiales des images.
\end{enumerate}
Ces mécanismes constituent les piliers des architectures modernes comme
\textbf{ResNet} ou les \textbf{Transformers}.
\end{resultbox}

\end{document}